%% Sample Exam 4 -- Medical Image Analysis (ETH Zurich, 2026)
%%
%% Written in the style of Example_Exam_with_Answers.pdf.
%%
%% Compile the version WITH solutions:      pdflatex "\def\withanswers{}%% Sample Exam 4 -- Medical Image Analysis (ETH Zurich, 2026)
%%
%% Written in the style of Example_Exam_with_Answers.pdf.
%%
%% Compile the version WITH solutions:      pdflatex "\def\withanswers{}%% Sample Exam 4 -- Medical Image Analysis (ETH Zurich, 2026)
%%
%% Written in the style of Example_Exam_with_Answers.pdf.
%%
%% Compile the version WITH solutions:      pdflatex "\def\withanswers{}%% Sample Exam 4 -- Medical Image Analysis (ETH Zurich, 2026)
%%
%% Written in the style of Example_Exam_with_Answers.pdf.
%%
%% Compile the version WITH solutions:      pdflatex "\def\withanswers{}\input{Sample_Exam_4.tex}"
%% Compile the blank version (no answers):  pdflatex Sample_Exam_4.tex
%%
\ifx\withanswers\undefined
  \documentclass[11pt]{exam}
\else
  \documentclass[11pt,answers]{exam}
\fi

\usepackage[margin=2.4cm]{geometry}
\usepackage{amsmath,amssymb}
\usepackage{array}
\usepackage[T1]{fontenc}

% ---------------------------------------------------------------- page style
\pagestyle{foot}
\firstpageheader{}{}{}
\firstpagefooter{}{}{}
\runningheader{}{}{}
\runningfooter{}{}{Page \thepage}

% -------------------------------------------------------------- multiplechoice
% Reproduce the "(A) (B) (C) ..." labelling of the original exam.
\renewcommand{\choicelabel}{$\bigcirc$~(\Alph{choice})}

% ---------------------------------------------------------- true/false macros
% \TFT -> the statement is TRUE, \TFF -> the statement is FALSE.
% In the answer version the correct word is set in capitals and bold.
\newcommand{\TFT}{\ifprintanswers\textbf{TRUE}\hspace{1.1em}False\else True\hspace{1.1em}False\fi}
\newcommand{\TFF}{True\hspace{1.1em}\ifprintanswers\textbf{FALSE}\else False\fi}

% ------------------------------------------------------------- writing space
% Reserves room to write in the blank version; collapses in the answer version.
\newcommand{\answerspace}[1]{\ifprintanswers\else\vspace{#1}\fi}

\begin{document}

% ------------------------------------------------------------------ title page
\begin{center}
  {\LARGE\bfseries Sample Exam}\\[2mm]
  {\Large Medical Image Analysis}\\[6mm]
\end{center}

\vspace{4mm}
\noindent Student name:\hrulefill

\vspace{6mm}
\noindent Immatriculation number:\hrulefill

\vspace{10mm}

\begin{center}
  {\large\bfseries Exam Questions}
\end{center}

\noindent\textbf{Note:} All the questions have equal weight.

\vspace{4mm}

\begin{questions}

% =============================================================== Question 1 ===
\question
In 2D acquisitions the slices are acquired one after the other in a pre-defined direction and
stacked into a 3D volume; in axial acquisitions the slices lie in the transverse plane. Voxel
sizes are given as in-plane (two values) and through-plane (one value, the slice thickness), and
the field-of-view is image size $\times$ voxel size. A \emph{dynamic} (4D) acquisition repeats the
whole volume a number of times, one volume per time point.

\vspace{1mm}
\noindent An axial dynamic acquisition is performed with in-plane voxel size
$1.5 \times 1.5$~mm$^2$ and 128 pixels in both directions, through-plane voxel size $5.0$~mm with
30 slices, and 200 time points acquired at a temporal resolution of $2.5$~s per volume.

\begin{parts}

  \part What is the field-of-view in the head-to-toe, left-to-right and anterior-to-posterior
  directions, respectively?

  \begin{solution}
    Answer: $150 \times 192 \times 192$~mm$^3$.

    The through-plane direction of an axial acquisition is head-to-toe: $30 \times 5.0 = 150$~mm.
    Both in-plane directions give $128 \times 1.5 = 192$~mm.
  \end{solution}
\answerspace{2.2cm}

  \part How many voxels does the complete 4D data set contain, and how long does the acquisition
  take?

  \begin{solution}
    Answer: $98\,304\,000$ voxels, and $500$~s $= 8$~min $20$~s.

    Per volume: $128 \times 128 \times 30 = 491\,520$ voxels. Over 200 time points:
    $491\,520 \times 200 = 98\,304\,000$. Duration: $200 \times 2.5 = 500$~s.
  \end{solution}
\answerspace{2.8cm}

  \part A second scan of the same patient is acquired at $1.0 \times 1.0 \times 1.0$~mm$^3$
  isotropic. To compare measurements between the two scans, which property of the images has to
  be matched? Give the resize factors needed to bring one volume of the dynamic acquisition onto
  that grid, and the resulting image size.

  \begin{solution}
    Answer: the \emph{voxel size} has to be matched (never the image size). Resize factors
    $5.0 \times 1.5 \times 1.5$ along head-to-toe, left-to-right and anterior-to-posterior,
    giving $150 \times 192 \times 192$~voxels.

    The resize factor is the ratio of the old voxel size to the new one: $5.0/1.0 = 5.0$
    head-to-toe and $1.5/1.0 = 1.5$ in the two in-plane directions. Applying these:
    $30 \cdot 5.0 = 150$ and $128 \cdot 1.5 = 192$ twice. The field-of-view is unchanged, as it
    must be --- only the sampling of it has changed.
  \end{solution}
\answerspace{4.2cm}

\end{parts}

% =============================================================== Question 2 ===
\question
Otsu's method selects a threshold for foreground--background separation automatically from the
image histogram. Let $p_1, p_2$ be the numbers of pixels assigned to the two groups by a candidate
threshold and $\sigma_1^2, \sigma_2^2$ the intensity variances within those groups. Which of the
two criteria below is the one Otsu's method uses? Please circle the correct choice.

\vspace{2mm}
\begin{center}
  $\displaystyle \min_\tau\ \ p_1 \sigma_1^2 + p_2 \sigma_2^2$
  \hspace{3.2cm}
  $\displaystyle \max_\tau\ \ p_1 \sigma_1^2 + p_2 \sigma_2^2$
\end{center}
\vspace{2mm}

\begin{solution}
  Answer: Left one --- Otsu's method \emph{minimises} the intra-class variance.

  A good threshold is one for which the two resulting groups are each as homogeneous as possible,
  i.e.\ the spread of intensities \emph{within} a group is small. Maximising the same quantity
  would deliberately choose the threshold that makes the two groups as internally inconsistent as
  possible, which is the worst threshold available.
\end{solution}
\answerspace{1.5cm}

% =============================================================== Question 3 ===
\question
Denoising can be written either as an energy minimisation or as a maximum-a-posteriori estimate.
$J$ is the observed noisy image and $I$ the image we want to recover.

\vspace{2mm}
\begin{center}
  $\displaystyle \max_I\ \log P(J \mid I) + \log P(I)$
  \hspace{2.0cm}
  $\displaystyle \min_I\ \lambda\, D(I,J) + R(I)$
\end{center}
\vspace{2mm}

\begin{parts}
  \part Please indicate which term corresponds to which term in the two formulations.

  \begin{solution}
    Answer:
    \begin{align*}
      \log P(J \mid I) &\ \Longleftrightarrow\ \lambda\, D(I,J) &&\text{(data consistency)}\\[1mm]
      \log P(I)        &\ \Longleftrightarrow\ R(I)             &&\text{(regularisation)}
    \end{align*}

    The likelihood encodes how the noise process works; the prior encodes which images are
    considered plausible before any data is seen.
  \end{solution}
\answerspace{2.4cm}

  \part For total variation denoising the data term is $D(I,J) = \|I - J\|_2^2$ and the
  corresponding likelihood is Gaussian, $P(J \mid I) = \mathcal{N}(J;\, I,\, \sigma^2)$. What is
  $\lambda$ in terms of $\sigma$, and what does that mean in words?

  \begin{solution}
    Answer: $\lambda = 1/(2\sigma^2)$.

    Taking $-\log$ of a Gaussian likelihood gives $\|I-J\|_2^2 / (2\sigma^2)$ up to a constant, so
    the weight in front of the data term \emph{is} the inverse noise variance. In words: the
    weighting between data term and regulariser is not an arbitrary tuning knob --- it states how
    much noise you believe is in the observation. A noisier image (large $\sigma$) means a small
    $\lambda$, which trusts the data less and lets the prior smooth more.
  \end{solution}
\answerspace{3.0cm}
\end{parts}

% =============================================================== Question 4 ===
\question
Please indicate whether the statement is true or false by circling the correct choice in each of
the following. $A$ denotes a binary image viewed as a set of pixels, $B$ a structural element, and
$B_x$ the structural element translated to $x$.

\begin{parts}
  \part \TFT\ \ The dilation of $A$ by $B$ is the set of all $x$ for which $B_x$ has a non-empty
  intersection with $A$.

  \part \TFF\ \ The opening of $A$ by $B$ is a dilation followed by an erosion, and it fills small
  gaps and holes.

  \part \TFF\ \ Morphological operations are linear and shift-invariant.

  \part \TFT\ \ Applying a median filter to a label map cannot introduce labels that were not
  already present in it.
\end{parts}

\begin{solution}
  (a) TRUE --- $A \oplus B \triangleq \{x \mid B_x \cap A \neq \emptyset\}$. Erosion is the
  stricter condition $A \ominus B \triangleq \{x \mid B_x \subseteq A\}$.

  (b) FALSE --- both halves are wrong, and they are wrong together. \emph{Opening} is an erosion
  followed by a dilation, $(A \ominus B) \oplus B$, and it removes small islands and protrusions.
  \emph{Closing} is the other order, $(A \oplus B) \ominus B$, and that is the one that fills small
  gaps and holes.

  (c) FALSE --- they are shift-invariant, but they are emphatically \emph{non-linear}. That is
  precisely why they can remove an isolated island without blurring the boundary of the structure
  next to it, which no linear filter can do.

  (d) TRUE --- the median of a set of values is always one of those values, so no new value can
  appear. This is what makes median filtering a safe post-processing step on a segmentation: it
  removes isolated mislabelled voxels without inventing a label that does not exist. Averaging
  filters do not have this property.
\end{solution}

% =============================================================== Question 5 ===
\question
In landmark-based registration a set of corresponding landmark pairs is used to determine the
transformation parameters. How many corresponding landmark pairs are needed for the system to be
exactly determined when the transformation model is (i) a 2D affine transformation, and (ii) a 3D
affine transformation? Explain the counting.

\begin{solution}
  Answer: (i) \textbf{3} landmark pairs in 2D, (ii) \textbf{4} landmark pairs in 3D.

  A 2D affine transformation has 6 free parameters ($A \in \mathbb{R}^{2\times2}$ plus
  $t \in \mathbb{R}^2$). Each landmark correspondence must match both coordinates of the point and
  therefore contributes 2 scalar equations. Setting $2N = 6$ gives $N = 3$.

  A 3D affine transformation has 12 free parameters ($A \in \mathbb{R}^{3\times3}$ plus
  $t \in \mathbb{R}^3$), and each correspondence contributes 3 equations, so $3N = 12$ gives
  $N = 4$.

  With more landmarks than this the system becomes over-determined and is solved in the
  least-squares sense: the landmarks themselves are then no longer aligned exactly, but the rest of
  the image is aligned better --- which is usually what you actually want.
\end{solution}
\answerspace{4.0cm}

% =============================================================== Question 6 ===
\question
Optimisation in image registration is difficult, and the objective function generally has local
minima that a gradient-based optimiser can get stuck in. Please describe the two practical
strategies presented in the lecture for dealing with this, and briefly explain why each of them
helps. You do not need to write down the equations.

\begin{solution}
  Answer: \emph{Sequential optimisation.} First align the images with a \textbf{linear}
  transformation model (rigid or affine), and only then run the deformable, non-linear
  registration. The linear model has very few parameters and removes the large global
  misalignment, so the non-linear stage starts close to the correct answer and only has to
  explain the local differences that remain. Running the deformable model straight away would ask
  it to absorb a global rotation or translation through local deformation, which is both harder to
  optimise and physically wrong.

  \emph{Multi-scale optimisation.} Build an image pyramid and start the registration at the
  \textbf{coarsest} scale, using the result of each scale to initialise the next finer one,
  $\theta^{(k-1)}_0 = \theta^{(k)}_*$. A coarse image has fewer details, so its objective function
  is smoother and has fewer local minima; the coarse solution captures the large-scale alignment
  and places the finer optimisation inside the basin of the correct minimum.

  The two are routinely combined, and both are really the same idea: solve an easier version of the
  problem first and use its answer as the starting point for the harder one.
\end{solution}
\answerspace{7.5cm}

% =============================================================== Question 7 ===
\question
Convolutional networks reduce the spatial dimensions of their feature maps either with pooling or
with strided convolutions. Which of the following statements is \textbf{inaccurate}?

\begin{choices}
  \choice Max pooling provides partial local translation invariance: the position of the maximum
  can move within the pooling window without changing the pooled value.
  \choice Pooling is applied to each channel separately and has no learnable parameters at all.
  \choice A strided convolution reduces the spatial dimensions and, like pooling, buys a degree of
  translation invariance.
  \choice Average pooling is a linear operation, whereas max pooling is not.
\end{choices}

\begin{solution}
  Answer: C.

  Strides help with dimensionality reduction, and that is all they do. Increasing the stride simply
  samples the convolution at fewer positions --- it discards information (the more so the larger the
  stride) and gives no invariance whatsoever, because shifting the input by one pixel changes which
  positions get sampled. Max pooling is different precisely because of the maximum: the largest
  activation inside the window can move around inside that window and the output does not change.
  That is a real, if partial and purely local, invariance.
\end{solution}

% =============================================================== Question 8 ===
\question
Three standard measures are used to judge how well a statistical shape model has embedded the
training data into a lower-dimensional space. Which of the following statements is
\textbf{incorrect}?

\begin{choices}
  \choice Compactness measures how much of the total variance is explained by the first $t$ modes,
  $C(t) = \sum_{i \leq t} \lambda_i / \lambda_{total}$; a compact model needs few parameters to
  generate new shape instances.
  \choice Generalisation measures the ability of the model to represent instances that were not in
  the training set, and is estimated with leave-one-out reconstruction.
  \choice Specificity measures how accurately the model reconstructs the training shapes
  themselves; a perfectly specific model reproduces every training instance exactly.
  \choice Principal component analysis assumes that the data follow a Gaussian distribution, and
  produces an orthogonal lower-dimensional coordinate system that maximises the variance along each
  successive axis.
\end{choices}

\begin{solution}
  Answer: C.

  Specificity asks the opposite question: does the model generate \emph{only valid} shapes? It is
  measured by drawing a large number of random instances from the model and, for each, computing
  the distance to the closest shape in the training set. A model that can generate anatomically
  impossible shapes scores badly, however well it reconstructs the training data.

  Note the tension this creates with (B): generalisation improves as you keep more modes,
  specificity degrades. The number of modes retained is a trade-off between the two.
\end{solution}

% =============================================================== Question 9 ===
\question
Which of the following statements about positional encoding and the structure of a transformer
layer is \textbf{incorrect}?

\begin{choices}
  \choice Without positional encoding a transformer is equivariant to permutations of its input
  tokens, so ``I bought an apple watch'' and ``watch an apple I bought'' cannot be distinguished.
  \choice An ideal positional encoding should give a unique representation for each position, be
  bounded, generalise to sequences longer than those seen in training, and encode the relative
  positions of tokens.
  \choice In the sinusoidal scheme the positional vector is concatenated to the token embedding,
  which doubles the dimension of the representation entering the first transformer layer.
  \choice A transformer layer consists of multi-head self-attention followed by a per-token MLP,
  each of the two wrapped in a residual connection and followed by layer normalisation.
\end{choices}

\begin{solution}
  Answer: C.

  The positional vector is \textbf{added}, not concatenated: $\tilde{x}_i = x_i + r_i$ with
  $r_i, x_i \in \mathbb{R}^{D \times 1}$. The dimension of the representation is therefore
  unchanged. This is why the encoding has to be bounded --- an unbounded positional vector added to
  the embedding would swamp the token's own content.
\end{solution}

% ============================================================== Question 10 ===
\question
The magnitude image of an MR acquisition is formed from the real and imaginary channels as
$J(x) = \sqrt{(I_r(x) + \epsilon)^2 + (I_i(x) + \eta)^2}$, where $\epsilon, \eta \sim
\mathcal{N}(0, \sigma)$ are the noise in the two channels. What kind of noise does the resulting
magnitude image have, and why does it matter?

\begin{choices}
  \choice Additive Gaussian noise, since $\epsilon$ and $\eta$ are both Gaussian and the operations
  applied to them are continuous.
  \choice Multiplicative speckle noise, of the same type as in ultrasound and optical coherence
  tomography.
  \choice Rician noise --- taking the magnitude makes the noise non-Gaussian and signal-dependent,
  so denoising methods derived from an additive Gaussian model are mis-specified, particularly in
  low-signal regions.
  \choice Poisson noise, of the same type as the quantum noise in low-dose CT.
  \choice No noise at all, because the two independent noise terms cancel in the magnitude
  operation.
\end{choices}

\begin{solution}
  Answer: C.

  The noise is Gaussian in the real and imaginary channels, but the magnitude operation is
  non-linear, and a non-linear function of Gaussian variables is not Gaussian. The result is the
  \textbf{Rician} distribution. The practical consequence is that the noise is signal-dependent:
  where the true signal is strong the Rician distribution is close to Gaussian and an additive
  model is a fair approximation, but in dark, low-signal regions it is strongly skewed and even
  biases the mean intensity upwards. This is why specialised Rician denoisers exist for MRI.
\end{solution}

% ============================================================== Question 11 ===
\question
Consider a classification network that takes grey-scale input images of size $60 \times 60$ with
one input channel. All convolutions are \emph{valid} (no padding, $P = 0$) with stride
$S = 1$, and each is followed by max pooling with $K = 2 \times 2$ and $S = 2$:

\begin{itemize}
  \item \textbf{Conv1}: 12 kernels with $K = 5 \times 5$ \quad\textbf{Pool1}: $2\times2$, $S = 2$
  \item \textbf{Conv2}: 24 kernels with $K = 5 \times 5$ \quad\textbf{Pool2}: $2\times2$, $S = 2$
  \item \textbf{Conv3}: 48 kernels with $K = 3 \times 3$ \quad\textbf{Pool3}: $2\times2$, $S = 2$
\end{itemize}

\noindent Recall that a valid convolution gives $N_l = N_{l-1} - k_1 + 1$. Circle True or False:

\begin{parts}
  \part \TFT\ \ The output of Pool3 has a spatial size of $5 \times 5$.

  \part \TFF\ \ Each neuron in Pool3 has a global receptive field covering the entire
  $60 \times 60$ input image.

  \part \TFT\ \ The total number of convolutional kernel weights (excluding the biases) is
  $5 \cdot 5 \cdot 1 \cdot 12 + 5 \cdot 5 \cdot 12 \cdot 24 + 3 \cdot 3 \cdot 24 \cdot 48
  = 17\,868$.

  \part \TFF\ \ Replacing the three max-pooling layers by $2 \times 2$ convolutions with stride 2
  would leave the total number of learnable parameters unchanged, because pooling and strided
  convolution both simply downsample.
\end{parts}

\begin{solution}
  (a) TRUE --- tracking the sizes: $60 \xrightarrow{\text{Conv1}} 56 \xrightarrow{\text{Pool1}} 28
  \xrightarrow{\text{Conv2}} 24 \xrightarrow{\text{Pool2}} 12 \xrightarrow{\text{Conv3}} 10
  \xrightarrow{\text{Pool3}} 5$.

  (b) FALSE --- it is $28 \times 28$, less than a quarter of the image area. Track the receptive
  field $R$ and jump $j$ from $R = 1$, $j = 1$ with $R \leftarrow R + (K-1)\cdot j$ and
  $j \leftarrow j \cdot S$: Conv1 $\rightarrow R = 5,\ j = 1$; Pool1 $\rightarrow R = 6,\ j = 2$;
  Conv2 $\rightarrow R = 14,\ j = 2$; Pool2 $\rightarrow R = 16,\ j = 4$;
  Conv3 $\rightarrow R = 24,\ j = 4$; Pool3 $\rightarrow R = 28,\ j = 8$. Output size and receptive
  field are different things --- a $5 \times 5$ output does not mean each of those 25 neurons sees
  everything. (The global stride, incidentally, is $8 \times 8$.)

  (c) TRUE --- $300 + 7\,200 + 10\,368 = 17\,868$. Pooling layers contribute nothing.

  (d) FALSE --- max pooling has \emph{no} learnable parameters, whereas convolutions do. The
  replacement would add
  $2 \cdot 2 \cdot 12 \cdot 12 + 2 \cdot 2 \cdot 24 \cdot 24 + 2 \cdot 2 \cdot 48 \cdot 48
  = 576 + 2\,304 + 9\,216 = 12\,096$ new weights, an increase of about 68\%. The two operations
  also differ in behaviour, not just in cost: pooling buys partial local translation invariance,
  strided convolution does not.
\end{solution}

% ============================================================== Question 12 ===
\question
Which of the following statements on domain shift, adaptation and uncertainty are correct (multiple
answers are possible)?

\begin{choices}
  \choice Acquisition shift is a change in $P_D(X \mid Z)$ --- scanner, resolution, contrast,
  protocol --- and is the kind of dataset shift encountered most often in medical imaging.
  \choice Prevalence shift is a change in $P_D(Y \mid X)$, caused for instance by differences in
  annotation policy or in the experience of the annotator.
  \choice The Expected Calibration Error measures the average gap between the confidence a model
  assigns to its predictions and the accuracy it actually achieves at that confidence.
  \choice Adapting only the batch-normalisation parameters of a network is not possible, because
  batch normalisation has no learnable parameters.
  \choice A deep ensemble approximates $p(\theta \mid D, M)$ by training the same architecture
  several times from different random initialisations, on the assumption that the different runs
  land in different modes of the parameter posterior.
\end{choices}

\begin{solution}
  Answer: (A), (C) and (E) are correct.

  (B) confuses two rows of the taxonomy. \emph{Prevalence} shift is a change in $P_D(Y)$ --- the
  case--control balance, or how the target population was selected. A change in $P_D(Y \mid X)$
  driven by annotation policy or annotator experience is \emph{annotation} shift.

  (D) is wrong: batch normalisation computes
  $\mathrm{BN}(a_l) = \gamma\,(a_l - \mu_l)/\sqrt{\sigma_l^2 + \epsilon} + \beta$, and the scale
  $\gamma$ and shift $\beta$ are learnable. Adapting only those is a standard lightweight strategy,
  attractive because it touches very few parameters and leaves the expensively trained weights
  alone --- though it is not obvious that they are the right parameters for correcting a change in
  intensity characteristics.
\end{solution}

\end{questions}

\end{document}
"
%% Compile the blank version (no answers):  pdflatex Sample_Exam_4.tex
%%
\ifx\withanswers\undefined
  \documentclass[11pt]{exam}
\else
  \documentclass[11pt,answers]{exam}
\fi

\usepackage[margin=2.4cm]{geometry}
\usepackage{amsmath,amssymb}
\usepackage{array}
\usepackage[T1]{fontenc}

% ---------------------------------------------------------------- page style
\pagestyle{foot}
\firstpageheader{}{}{}
\firstpagefooter{}{}{}
\runningheader{}{}{}
\runningfooter{}{}{Page \thepage}

% -------------------------------------------------------------- multiplechoice
% Reproduce the "(A) (B) (C) ..." labelling of the original exam.
\renewcommand{\choicelabel}{$\bigcirc$~(\Alph{choice})}

% ---------------------------------------------------------- true/false macros
% \TFT -> the statement is TRUE, \TFF -> the statement is FALSE.
% In the answer version the correct word is set in capitals and bold.
\newcommand{\TFT}{\ifprintanswers\textbf{TRUE}\hspace{1.1em}False\else True\hspace{1.1em}False\fi}
\newcommand{\TFF}{True\hspace{1.1em}\ifprintanswers\textbf{FALSE}\else False\fi}

% ------------------------------------------------------------- writing space
% Reserves room to write in the blank version; collapses in the answer version.
\newcommand{\answerspace}[1]{\ifprintanswers\else\vspace{#1}\fi}

\begin{document}

% ------------------------------------------------------------------ title page
\begin{center}
  {\LARGE\bfseries Sample Exam}\\[2mm]
  {\Large Medical Image Analysis}\\[6mm]
\end{center}

\vspace{4mm}
\noindent Student name:\hrulefill

\vspace{6mm}
\noindent Immatriculation number:\hrulefill

\vspace{10mm}

\begin{center}
  {\large\bfseries Exam Questions}
\end{center}

\noindent\textbf{Note:} All the questions have equal weight.

\vspace{4mm}

\begin{questions}

% =============================================================== Question 1 ===
\question
In 2D acquisitions the slices are acquired one after the other in a pre-defined direction and
stacked into a 3D volume; in axial acquisitions the slices lie in the transverse plane. Voxel
sizes are given as in-plane (two values) and through-plane (one value, the slice thickness), and
the field-of-view is image size $\times$ voxel size. A \emph{dynamic} (4D) acquisition repeats the
whole volume a number of times, one volume per time point.

\vspace{1mm}
\noindent An axial dynamic acquisition is performed with in-plane voxel size
$1.5 \times 1.5$~mm$^2$ and 128 pixels in both directions, through-plane voxel size $5.0$~mm with
30 slices, and 200 time points acquired at a temporal resolution of $2.5$~s per volume.

\begin{parts}

  \part What is the field-of-view in the head-to-toe, left-to-right and anterior-to-posterior
  directions, respectively?

  \begin{solution}
    Answer: $150 \times 192 \times 192$~mm$^3$.

    The through-plane direction of an axial acquisition is head-to-toe: $30 \times 5.0 = 150$~mm.
    Both in-plane directions give $128 \times 1.5 = 192$~mm.
  \end{solution}
\answerspace{2.2cm}

  \part How many voxels does the complete 4D data set contain, and how long does the acquisition
  take?

  \begin{solution}
    Answer: $98\,304\,000$ voxels, and $500$~s $= 8$~min $20$~s.

    Per volume: $128 \times 128 \times 30 = 491\,520$ voxels. Over 200 time points:
    $491\,520 \times 200 = 98\,304\,000$. Duration: $200 \times 2.5 = 500$~s.
  \end{solution}
\answerspace{2.8cm}

  \part A second scan of the same patient is acquired at $1.0 \times 1.0 \times 1.0$~mm$^3$
  isotropic. To compare measurements between the two scans, which property of the images has to
  be matched? Give the resize factors needed to bring one volume of the dynamic acquisition onto
  that grid, and the resulting image size.

  \begin{solution}
    Answer: the \emph{voxel size} has to be matched (never the image size). Resize factors
    $5.0 \times 1.5 \times 1.5$ along head-to-toe, left-to-right and anterior-to-posterior,
    giving $150 \times 192 \times 192$~voxels.

    The resize factor is the ratio of the old voxel size to the new one: $5.0/1.0 = 5.0$
    head-to-toe and $1.5/1.0 = 1.5$ in the two in-plane directions. Applying these:
    $30 \cdot 5.0 = 150$ and $128 \cdot 1.5 = 192$ twice. The field-of-view is unchanged, as it
    must be --- only the sampling of it has changed.
  \end{solution}
\answerspace{4.2cm}

\end{parts}

% =============================================================== Question 2 ===
\question
Otsu's method selects a threshold for foreground--background separation automatically from the
image histogram. Let $p_1, p_2$ be the numbers of pixels assigned to the two groups by a candidate
threshold and $\sigma_1^2, \sigma_2^2$ the intensity variances within those groups. Which of the
two criteria below is the one Otsu's method uses? Please circle the correct choice.

\vspace{2mm}
\begin{center}
  $\displaystyle \min_\tau\ \ p_1 \sigma_1^2 + p_2 \sigma_2^2$
  \hspace{3.2cm}
  $\displaystyle \max_\tau\ \ p_1 \sigma_1^2 + p_2 \sigma_2^2$
\end{center}
\vspace{2mm}

\begin{solution}
  Answer: Left one --- Otsu's method \emph{minimises} the intra-class variance.

  A good threshold is one for which the two resulting groups are each as homogeneous as possible,
  i.e.\ the spread of intensities \emph{within} a group is small. Maximising the same quantity
  would deliberately choose the threshold that makes the two groups as internally inconsistent as
  possible, which is the worst threshold available.
\end{solution}
\answerspace{1.5cm}

% =============================================================== Question 3 ===
\question
Denoising can be written either as an energy minimisation or as a maximum-a-posteriori estimate.
$J$ is the observed noisy image and $I$ the image we want to recover.

\vspace{2mm}
\begin{center}
  $\displaystyle \max_I\ \log P(J \mid I) + \log P(I)$
  \hspace{2.0cm}
  $\displaystyle \min_I\ \lambda\, D(I,J) + R(I)$
\end{center}
\vspace{2mm}

\begin{parts}
  \part Please indicate which term corresponds to which term in the two formulations.

  \begin{solution}
    Answer:
    \begin{align*}
      \log P(J \mid I) &\ \Longleftrightarrow\ \lambda\, D(I,J) &&\text{(data consistency)}\\[1mm]
      \log P(I)        &\ \Longleftrightarrow\ R(I)             &&\text{(regularisation)}
    \end{align*}

    The likelihood encodes how the noise process works; the prior encodes which images are
    considered plausible before any data is seen.
  \end{solution}
\answerspace{2.4cm}

  \part For total variation denoising the data term is $D(I,J) = \|I - J\|_2^2$ and the
  corresponding likelihood is Gaussian, $P(J \mid I) = \mathcal{N}(J;\, I,\, \sigma^2)$. What is
  $\lambda$ in terms of $\sigma$, and what does that mean in words?

  \begin{solution}
    Answer: $\lambda = 1/(2\sigma^2)$.

    Taking $-\log$ of a Gaussian likelihood gives $\|I-J\|_2^2 / (2\sigma^2)$ up to a constant, so
    the weight in front of the data term \emph{is} the inverse noise variance. In words: the
    weighting between data term and regulariser is not an arbitrary tuning knob --- it states how
    much noise you believe is in the observation. A noisier image (large $\sigma$) means a small
    $\lambda$, which trusts the data less and lets the prior smooth more.
  \end{solution}
\answerspace{3.0cm}
\end{parts}

% =============================================================== Question 4 ===
\question
Please indicate whether the statement is true or false by circling the correct choice in each of
the following. $A$ denotes a binary image viewed as a set of pixels, $B$ a structural element, and
$B_x$ the structural element translated to $x$.

\begin{parts}
  \part \TFT\ \ The dilation of $A$ by $B$ is the set of all $x$ for which $B_x$ has a non-empty
  intersection with $A$.

  \part \TFF\ \ The opening of $A$ by $B$ is a dilation followed by an erosion, and it fills small
  gaps and holes.

  \part \TFF\ \ Morphological operations are linear and shift-invariant.

  \part \TFT\ \ Applying a median filter to a label map cannot introduce labels that were not
  already present in it.
\end{parts}

\begin{solution}
  (a) TRUE --- $A \oplus B \triangleq \{x \mid B_x \cap A \neq \emptyset\}$. Erosion is the
  stricter condition $A \ominus B \triangleq \{x \mid B_x \subseteq A\}$.

  (b) FALSE --- both halves are wrong, and they are wrong together. \emph{Opening} is an erosion
  followed by a dilation, $(A \ominus B) \oplus B$, and it removes small islands and protrusions.
  \emph{Closing} is the other order, $(A \oplus B) \ominus B$, and that is the one that fills small
  gaps and holes.

  (c) FALSE --- they are shift-invariant, but they are emphatically \emph{non-linear}. That is
  precisely why they can remove an isolated island without blurring the boundary of the structure
  next to it, which no linear filter can do.

  (d) TRUE --- the median of a set of values is always one of those values, so no new value can
  appear. This is what makes median filtering a safe post-processing step on a segmentation: it
  removes isolated mislabelled voxels without inventing a label that does not exist. Averaging
  filters do not have this property.
\end{solution}

% =============================================================== Question 5 ===
\question
In landmark-based registration a set of corresponding landmark pairs is used to determine the
transformation parameters. How many corresponding landmark pairs are needed for the system to be
exactly determined when the transformation model is (i) a 2D affine transformation, and (ii) a 3D
affine transformation? Explain the counting.

\begin{solution}
  Answer: (i) \textbf{3} landmark pairs in 2D, (ii) \textbf{4} landmark pairs in 3D.

  A 2D affine transformation has 6 free parameters ($A \in \mathbb{R}^{2\times2}$ plus
  $t \in \mathbb{R}^2$). Each landmark correspondence must match both coordinates of the point and
  therefore contributes 2 scalar equations. Setting $2N = 6$ gives $N = 3$.

  A 3D affine transformation has 12 free parameters ($A \in \mathbb{R}^{3\times3}$ plus
  $t \in \mathbb{R}^3$), and each correspondence contributes 3 equations, so $3N = 12$ gives
  $N = 4$.

  With more landmarks than this the system becomes over-determined and is solved in the
  least-squares sense: the landmarks themselves are then no longer aligned exactly, but the rest of
  the image is aligned better --- which is usually what you actually want.
\end{solution}
\answerspace{4.0cm}

% =============================================================== Question 6 ===
\question
Optimisation in image registration is difficult, and the objective function generally has local
minima that a gradient-based optimiser can get stuck in. Please describe the two practical
strategies presented in the lecture for dealing with this, and briefly explain why each of them
helps. You do not need to write down the equations.

\begin{solution}
  Answer: \emph{Sequential optimisation.} First align the images with a \textbf{linear}
  transformation model (rigid or affine), and only then run the deformable, non-linear
  registration. The linear model has very few parameters and removes the large global
  misalignment, so the non-linear stage starts close to the correct answer and only has to
  explain the local differences that remain. Running the deformable model straight away would ask
  it to absorb a global rotation or translation through local deformation, which is both harder to
  optimise and physically wrong.

  \emph{Multi-scale optimisation.} Build an image pyramid and start the registration at the
  \textbf{coarsest} scale, using the result of each scale to initialise the next finer one,
  $\theta^{(k-1)}_0 = \theta^{(k)}_*$. A coarse image has fewer details, so its objective function
  is smoother and has fewer local minima; the coarse solution captures the large-scale alignment
  and places the finer optimisation inside the basin of the correct minimum.

  The two are routinely combined, and both are really the same idea: solve an easier version of the
  problem first and use its answer as the starting point for the harder one.
\end{solution}
\answerspace{7.5cm}

% =============================================================== Question 7 ===
\question
Convolutional networks reduce the spatial dimensions of their feature maps either with pooling or
with strided convolutions. Which of the following statements is \textbf{inaccurate}?

\begin{choices}
  \choice Max pooling provides partial local translation invariance: the position of the maximum
  can move within the pooling window without changing the pooled value.
  \choice Pooling is applied to each channel separately and has no learnable parameters at all.
  \choice A strided convolution reduces the spatial dimensions and, like pooling, buys a degree of
  translation invariance.
  \choice Average pooling is a linear operation, whereas max pooling is not.
\end{choices}

\begin{solution}
  Answer: C.

  Strides help with dimensionality reduction, and that is all they do. Increasing the stride simply
  samples the convolution at fewer positions --- it discards information (the more so the larger the
  stride) and gives no invariance whatsoever, because shifting the input by one pixel changes which
  positions get sampled. Max pooling is different precisely because of the maximum: the largest
  activation inside the window can move around inside that window and the output does not change.
  That is a real, if partial and purely local, invariance.
\end{solution}

% =============================================================== Question 8 ===
\question
Three standard measures are used to judge how well a statistical shape model has embedded the
training data into a lower-dimensional space. Which of the following statements is
\textbf{incorrect}?

\begin{choices}
  \choice Compactness measures how much of the total variance is explained by the first $t$ modes,
  $C(t) = \sum_{i \leq t} \lambda_i / \lambda_{total}$; a compact model needs few parameters to
  generate new shape instances.
  \choice Generalisation measures the ability of the model to represent instances that were not in
  the training set, and is estimated with leave-one-out reconstruction.
  \choice Specificity measures how accurately the model reconstructs the training shapes
  themselves; a perfectly specific model reproduces every training instance exactly.
  \choice Principal component analysis assumes that the data follow a Gaussian distribution, and
  produces an orthogonal lower-dimensional coordinate system that maximises the variance along each
  successive axis.
\end{choices}

\begin{solution}
  Answer: C.

  Specificity asks the opposite question: does the model generate \emph{only valid} shapes? It is
  measured by drawing a large number of random instances from the model and, for each, computing
  the distance to the closest shape in the training set. A model that can generate anatomically
  impossible shapes scores badly, however well it reconstructs the training data.

  Note the tension this creates with (B): generalisation improves as you keep more modes,
  specificity degrades. The number of modes retained is a trade-off between the two.
\end{solution}

% =============================================================== Question 9 ===
\question
Which of the following statements about positional encoding and the structure of a transformer
layer is \textbf{incorrect}?

\begin{choices}
  \choice Without positional encoding a transformer is equivariant to permutations of its input
  tokens, so ``I bought an apple watch'' and ``watch an apple I bought'' cannot be distinguished.
  \choice An ideal positional encoding should give a unique representation for each position, be
  bounded, generalise to sequences longer than those seen in training, and encode the relative
  positions of tokens.
  \choice In the sinusoidal scheme the positional vector is concatenated to the token embedding,
  which doubles the dimension of the representation entering the first transformer layer.
  \choice A transformer layer consists of multi-head self-attention followed by a per-token MLP,
  each of the two wrapped in a residual connection and followed by layer normalisation.
\end{choices}

\begin{solution}
  Answer: C.

  The positional vector is \textbf{added}, not concatenated: $\tilde{x}_i = x_i + r_i$ with
  $r_i, x_i \in \mathbb{R}^{D \times 1}$. The dimension of the representation is therefore
  unchanged. This is why the encoding has to be bounded --- an unbounded positional vector added to
  the embedding would swamp the token's own content.
\end{solution}

% ============================================================== Question 10 ===
\question
The magnitude image of an MR acquisition is formed from the real and imaginary channels as
$J(x) = \sqrt{(I_r(x) + \epsilon)^2 + (I_i(x) + \eta)^2}$, where $\epsilon, \eta \sim
\mathcal{N}(0, \sigma)$ are the noise in the two channels. What kind of noise does the resulting
magnitude image have, and why does it matter?

\begin{choices}
  \choice Additive Gaussian noise, since $\epsilon$ and $\eta$ are both Gaussian and the operations
  applied to them are continuous.
  \choice Multiplicative speckle noise, of the same type as in ultrasound and optical coherence
  tomography.
  \choice Rician noise --- taking the magnitude makes the noise non-Gaussian and signal-dependent,
  so denoising methods derived from an additive Gaussian model are mis-specified, particularly in
  low-signal regions.
  \choice Poisson noise, of the same type as the quantum noise in low-dose CT.
  \choice No noise at all, because the two independent noise terms cancel in the magnitude
  operation.
\end{choices}

\begin{solution}
  Answer: C.

  The noise is Gaussian in the real and imaginary channels, but the magnitude operation is
  non-linear, and a non-linear function of Gaussian variables is not Gaussian. The result is the
  \textbf{Rician} distribution. The practical consequence is that the noise is signal-dependent:
  where the true signal is strong the Rician distribution is close to Gaussian and an additive
  model is a fair approximation, but in dark, low-signal regions it is strongly skewed and even
  biases the mean intensity upwards. This is why specialised Rician denoisers exist for MRI.
\end{solution}

% ============================================================== Question 11 ===
\question
Consider a classification network that takes grey-scale input images of size $60 \times 60$ with
one input channel. All convolutions are \emph{valid} (no padding, $P = 0$) with stride
$S = 1$, and each is followed by max pooling with $K = 2 \times 2$ and $S = 2$:

\begin{itemize}
  \item \textbf{Conv1}: 12 kernels with $K = 5 \times 5$ \quad\textbf{Pool1}: $2\times2$, $S = 2$
  \item \textbf{Conv2}: 24 kernels with $K = 5 \times 5$ \quad\textbf{Pool2}: $2\times2$, $S = 2$
  \item \textbf{Conv3}: 48 kernels with $K = 3 \times 3$ \quad\textbf{Pool3}: $2\times2$, $S = 2$
\end{itemize}

\noindent Recall that a valid convolution gives $N_l = N_{l-1} - k_1 + 1$. Circle True or False:

\begin{parts}
  \part \TFT\ \ The output of Pool3 has a spatial size of $5 \times 5$.

  \part \TFF\ \ Each neuron in Pool3 has a global receptive field covering the entire
  $60 \times 60$ input image.

  \part \TFT\ \ The total number of convolutional kernel weights (excluding the biases) is
  $5 \cdot 5 \cdot 1 \cdot 12 + 5 \cdot 5 \cdot 12 \cdot 24 + 3 \cdot 3 \cdot 24 \cdot 48
  = 17\,868$.

  \part \TFF\ \ Replacing the three max-pooling layers by $2 \times 2$ convolutions with stride 2
  would leave the total number of learnable parameters unchanged, because pooling and strided
  convolution both simply downsample.
\end{parts}

\begin{solution}
  (a) TRUE --- tracking the sizes: $60 \xrightarrow{\text{Conv1}} 56 \xrightarrow{\text{Pool1}} 28
  \xrightarrow{\text{Conv2}} 24 \xrightarrow{\text{Pool2}} 12 \xrightarrow{\text{Conv3}} 10
  \xrightarrow{\text{Pool3}} 5$.

  (b) FALSE --- it is $28 \times 28$, less than a quarter of the image area. Track the receptive
  field $R$ and jump $j$ from $R = 1$, $j = 1$ with $R \leftarrow R + (K-1)\cdot j$ and
  $j \leftarrow j \cdot S$: Conv1 $\rightarrow R = 5,\ j = 1$; Pool1 $\rightarrow R = 6,\ j = 2$;
  Conv2 $\rightarrow R = 14,\ j = 2$; Pool2 $\rightarrow R = 16,\ j = 4$;
  Conv3 $\rightarrow R = 24,\ j = 4$; Pool3 $\rightarrow R = 28,\ j = 8$. Output size and receptive
  field are different things --- a $5 \times 5$ output does not mean each of those 25 neurons sees
  everything. (The global stride, incidentally, is $8 \times 8$.)

  (c) TRUE --- $300 + 7\,200 + 10\,368 = 17\,868$. Pooling layers contribute nothing.

  (d) FALSE --- max pooling has \emph{no} learnable parameters, whereas convolutions do. The
  replacement would add
  $2 \cdot 2 \cdot 12 \cdot 12 + 2 \cdot 2 \cdot 24 \cdot 24 + 2 \cdot 2 \cdot 48 \cdot 48
  = 576 + 2\,304 + 9\,216 = 12\,096$ new weights, an increase of about 68\%. The two operations
  also differ in behaviour, not just in cost: pooling buys partial local translation invariance,
  strided convolution does not.
\end{solution}

% ============================================================== Question 12 ===
\question
Which of the following statements on domain shift, adaptation and uncertainty are correct (multiple
answers are possible)?

\begin{choices}
  \choice Acquisition shift is a change in $P_D(X \mid Z)$ --- scanner, resolution, contrast,
  protocol --- and is the kind of dataset shift encountered most often in medical imaging.
  \choice Prevalence shift is a change in $P_D(Y \mid X)$, caused for instance by differences in
  annotation policy or in the experience of the annotator.
  \choice The Expected Calibration Error measures the average gap between the confidence a model
  assigns to its predictions and the accuracy it actually achieves at that confidence.
  \choice Adapting only the batch-normalisation parameters of a network is not possible, because
  batch normalisation has no learnable parameters.
  \choice A deep ensemble approximates $p(\theta \mid D, M)$ by training the same architecture
  several times from different random initialisations, on the assumption that the different runs
  land in different modes of the parameter posterior.
\end{choices}

\begin{solution}
  Answer: (A), (C) and (E) are correct.

  (B) confuses two rows of the taxonomy. \emph{Prevalence} shift is a change in $P_D(Y)$ --- the
  case--control balance, or how the target population was selected. A change in $P_D(Y \mid X)$
  driven by annotation policy or annotator experience is \emph{annotation} shift.

  (D) is wrong: batch normalisation computes
  $\mathrm{BN}(a_l) = \gamma\,(a_l - \mu_l)/\sqrt{\sigma_l^2 + \epsilon} + \beta$, and the scale
  $\gamma$ and shift $\beta$ are learnable. Adapting only those is a standard lightweight strategy,
  attractive because it touches very few parameters and leaves the expensively trained weights
  alone --- though it is not obvious that they are the right parameters for correcting a change in
  intensity characteristics.
\end{solution}

\end{questions}

\end{document}
"
%% Compile the blank version (no answers):  pdflatex Sample_Exam_4.tex
%%
\ifx\withanswers\undefined
  \documentclass[11pt]{exam}
\else
  \documentclass[11pt,answers]{exam}
\fi

\usepackage[margin=2.4cm]{geometry}
\usepackage{amsmath,amssymb}
\usepackage{array}
\usepackage[T1]{fontenc}

% ---------------------------------------------------------------- page style
\pagestyle{foot}
\firstpageheader{}{}{}
\firstpagefooter{}{}{}
\runningheader{}{}{}
\runningfooter{}{}{Page \thepage}

% -------------------------------------------------------------- multiplechoice
% Reproduce the "(A) (B) (C) ..." labelling of the original exam.
\renewcommand{\choicelabel}{$\bigcirc$~(\Alph{choice})}

% ---------------------------------------------------------- true/false macros
% \TFT -> the statement is TRUE, \TFF -> the statement is FALSE.
% In the answer version the correct word is set in capitals and bold.
\newcommand{\TFT}{\ifprintanswers\textbf{TRUE}\hspace{1.1em}False\else True\hspace{1.1em}False\fi}
\newcommand{\TFF}{True\hspace{1.1em}\ifprintanswers\textbf{FALSE}\else False\fi}

% ------------------------------------------------------------- writing space
% Reserves room to write in the blank version; collapses in the answer version.
\newcommand{\answerspace}[1]{\ifprintanswers\else\vspace{#1}\fi}

\begin{document}

% ------------------------------------------------------------------ title page
\begin{center}
  {\LARGE\bfseries Sample Exam}\\[2mm]
  {\Large Medical Image Analysis}\\[6mm]
\end{center}

\vspace{4mm}
\noindent Student name:\hrulefill

\vspace{6mm}
\noindent Immatriculation number:\hrulefill

\vspace{10mm}

\begin{center}
  {\large\bfseries Exam Questions}
\end{center}

\noindent\textbf{Note:} All the questions have equal weight.

\vspace{4mm}

\begin{questions}

% =============================================================== Question 1 ===
\question
In 2D acquisitions the slices are acquired one after the other in a pre-defined direction and
stacked into a 3D volume; in axial acquisitions the slices lie in the transverse plane. Voxel
sizes are given as in-plane (two values) and through-plane (one value, the slice thickness), and
the field-of-view is image size $\times$ voxel size. A \emph{dynamic} (4D) acquisition repeats the
whole volume a number of times, one volume per time point.

\vspace{1mm}
\noindent An axial dynamic acquisition is performed with in-plane voxel size
$1.5 \times 1.5$~mm$^2$ and 128 pixels in both directions, through-plane voxel size $5.0$~mm with
30 slices, and 200 time points acquired at a temporal resolution of $2.5$~s per volume.

\begin{parts}

  \part What is the field-of-view in the head-to-toe, left-to-right and anterior-to-posterior
  directions, respectively?

  \begin{solution}
    Answer: $150 \times 192 \times 192$~mm$^3$.

    The through-plane direction of an axial acquisition is head-to-toe: $30 \times 5.0 = 150$~mm.
    Both in-plane directions give $128 \times 1.5 = 192$~mm.
  \end{solution}
\answerspace{2.2cm}

  \part How many voxels does the complete 4D data set contain, and how long does the acquisition
  take?

  \begin{solution}
    Answer: $98\,304\,000$ voxels, and $500$~s $= 8$~min $20$~s.

    Per volume: $128 \times 128 \times 30 = 491\,520$ voxels. Over 200 time points:
    $491\,520 \times 200 = 98\,304\,000$. Duration: $200 \times 2.5 = 500$~s.
  \end{solution}
\answerspace{2.8cm}

  \part A second scan of the same patient is acquired at $1.0 \times 1.0 \times 1.0$~mm$^3$
  isotropic. To compare measurements between the two scans, which property of the images has to
  be matched? Give the resize factors needed to bring one volume of the dynamic acquisition onto
  that grid, and the resulting image size.

  \begin{solution}
    Answer: the \emph{voxel size} has to be matched (never the image size). Resize factors
    $5.0 \times 1.5 \times 1.5$ along head-to-toe, left-to-right and anterior-to-posterior,
    giving $150 \times 192 \times 192$~voxels.

    The resize factor is the ratio of the old voxel size to the new one: $5.0/1.0 = 5.0$
    head-to-toe and $1.5/1.0 = 1.5$ in the two in-plane directions. Applying these:
    $30 \cdot 5.0 = 150$ and $128 \cdot 1.5 = 192$ twice. The field-of-view is unchanged, as it
    must be --- only the sampling of it has changed.
  \end{solution}
\answerspace{4.2cm}

\end{parts}

% =============================================================== Question 2 ===
\question
Otsu's method selects a threshold for foreground--background separation automatically from the
image histogram. Let $p_1, p_2$ be the numbers of pixels assigned to the two groups by a candidate
threshold and $\sigma_1^2, \sigma_2^2$ the intensity variances within those groups. Which of the
two criteria below is the one Otsu's method uses? Please circle the correct choice.

\vspace{2mm}
\begin{center}
  $\displaystyle \min_\tau\ \ p_1 \sigma_1^2 + p_2 \sigma_2^2$
  \hspace{3.2cm}
  $\displaystyle \max_\tau\ \ p_1 \sigma_1^2 + p_2 \sigma_2^2$
\end{center}
\vspace{2mm}

\begin{solution}
  Answer: Left one --- Otsu's method \emph{minimises} the intra-class variance.

  A good threshold is one for which the two resulting groups are each as homogeneous as possible,
  i.e.\ the spread of intensities \emph{within} a group is small. Maximising the same quantity
  would deliberately choose the threshold that makes the two groups as internally inconsistent as
  possible, which is the worst threshold available.
\end{solution}
\answerspace{1.5cm}

% =============================================================== Question 3 ===
\question
Denoising can be written either as an energy minimisation or as a maximum-a-posteriori estimate.
$J$ is the observed noisy image and $I$ the image we want to recover.

\vspace{2mm}
\begin{center}
  $\displaystyle \max_I\ \log P(J \mid I) + \log P(I)$
  \hspace{2.0cm}
  $\displaystyle \min_I\ \lambda\, D(I,J) + R(I)$
\end{center}
\vspace{2mm}

\begin{parts}
  \part Please indicate which term corresponds to which term in the two formulations.

  \begin{solution}
    Answer:
    \begin{align*}
      \log P(J \mid I) &\ \Longleftrightarrow\ \lambda\, D(I,J) &&\text{(data consistency)}\\[1mm]
      \log P(I)        &\ \Longleftrightarrow\ R(I)             &&\text{(regularisation)}
    \end{align*}

    The likelihood encodes how the noise process works; the prior encodes which images are
    considered plausible before any data is seen.
  \end{solution}
\answerspace{2.4cm}

  \part For total variation denoising the data term is $D(I,J) = \|I - J\|_2^2$ and the
  corresponding likelihood is Gaussian, $P(J \mid I) = \mathcal{N}(J;\, I,\, \sigma^2)$. What is
  $\lambda$ in terms of $\sigma$, and what does that mean in words?

  \begin{solution}
    Answer: $\lambda = 1/(2\sigma^2)$.

    Taking $-\log$ of a Gaussian likelihood gives $\|I-J\|_2^2 / (2\sigma^2)$ up to a constant, so
    the weight in front of the data term \emph{is} the inverse noise variance. In words: the
    weighting between data term and regulariser is not an arbitrary tuning knob --- it states how
    much noise you believe is in the observation. A noisier image (large $\sigma$) means a small
    $\lambda$, which trusts the data less and lets the prior smooth more.
  \end{solution}
\answerspace{3.0cm}
\end{parts}

% =============================================================== Question 4 ===
\question
Please indicate whether the statement is true or false by circling the correct choice in each of
the following. $A$ denotes a binary image viewed as a set of pixels, $B$ a structural element, and
$B_x$ the structural element translated to $x$.

\begin{parts}
  \part \TFT\ \ The dilation of $A$ by $B$ is the set of all $x$ for which $B_x$ has a non-empty
  intersection with $A$.

  \part \TFF\ \ The opening of $A$ by $B$ is a dilation followed by an erosion, and it fills small
  gaps and holes.

  \part \TFF\ \ Morphological operations are linear and shift-invariant.

  \part \TFT\ \ Applying a median filter to a label map cannot introduce labels that were not
  already present in it.
\end{parts}

\begin{solution}
  (a) TRUE --- $A \oplus B \triangleq \{x \mid B_x \cap A \neq \emptyset\}$. Erosion is the
  stricter condition $A \ominus B \triangleq \{x \mid B_x \subseteq A\}$.

  (b) FALSE --- both halves are wrong, and they are wrong together. \emph{Opening} is an erosion
  followed by a dilation, $(A \ominus B) \oplus B$, and it removes small islands and protrusions.
  \emph{Closing} is the other order, $(A \oplus B) \ominus B$, and that is the one that fills small
  gaps and holes.

  (c) FALSE --- they are shift-invariant, but they are emphatically \emph{non-linear}. That is
  precisely why they can remove an isolated island without blurring the boundary of the structure
  next to it, which no linear filter can do.

  (d) TRUE --- the median of a set of values is always one of those values, so no new value can
  appear. This is what makes median filtering a safe post-processing step on a segmentation: it
  removes isolated mislabelled voxels without inventing a label that does not exist. Averaging
  filters do not have this property.
\end{solution}

% =============================================================== Question 5 ===
\question
In landmark-based registration a set of corresponding landmark pairs is used to determine the
transformation parameters. How many corresponding landmark pairs are needed for the system to be
exactly determined when the transformation model is (i) a 2D affine transformation, and (ii) a 3D
affine transformation? Explain the counting.

\begin{solution}
  Answer: (i) \textbf{3} landmark pairs in 2D, (ii) \textbf{4} landmark pairs in 3D.

  A 2D affine transformation has 6 free parameters ($A \in \mathbb{R}^{2\times2}$ plus
  $t \in \mathbb{R}^2$). Each landmark correspondence must match both coordinates of the point and
  therefore contributes 2 scalar equations. Setting $2N = 6$ gives $N = 3$.

  A 3D affine transformation has 12 free parameters ($A \in \mathbb{R}^{3\times3}$ plus
  $t \in \mathbb{R}^3$), and each correspondence contributes 3 equations, so $3N = 12$ gives
  $N = 4$.

  With more landmarks than this the system becomes over-determined and is solved in the
  least-squares sense: the landmarks themselves are then no longer aligned exactly, but the rest of
  the image is aligned better --- which is usually what you actually want.
\end{solution}
\answerspace{4.0cm}

% =============================================================== Question 6 ===
\question
Optimisation in image registration is difficult, and the objective function generally has local
minima that a gradient-based optimiser can get stuck in. Please describe the two practical
strategies presented in the lecture for dealing with this, and briefly explain why each of them
helps. You do not need to write down the equations.

\begin{solution}
  Answer: \emph{Sequential optimisation.} First align the images with a \textbf{linear}
  transformation model (rigid or affine), and only then run the deformable, non-linear
  registration. The linear model has very few parameters and removes the large global
  misalignment, so the non-linear stage starts close to the correct answer and only has to
  explain the local differences that remain. Running the deformable model straight away would ask
  it to absorb a global rotation or translation through local deformation, which is both harder to
  optimise and physically wrong.

  \emph{Multi-scale optimisation.} Build an image pyramid and start the registration at the
  \textbf{coarsest} scale, using the result of each scale to initialise the next finer one,
  $\theta^{(k-1)}_0 = \theta^{(k)}_*$. A coarse image has fewer details, so its objective function
  is smoother and has fewer local minima; the coarse solution captures the large-scale alignment
  and places the finer optimisation inside the basin of the correct minimum.

  The two are routinely combined, and both are really the same idea: solve an easier version of the
  problem first and use its answer as the starting point for the harder one.
\end{solution}
\answerspace{7.5cm}

% =============================================================== Question 7 ===
\question
Convolutional networks reduce the spatial dimensions of their feature maps either with pooling or
with strided convolutions. Which of the following statements is \textbf{inaccurate}?

\begin{choices}
  \choice Max pooling provides partial local translation invariance: the position of the maximum
  can move within the pooling window without changing the pooled value.
  \choice Pooling is applied to each channel separately and has no learnable parameters at all.
  \choice A strided convolution reduces the spatial dimensions and, like pooling, buys a degree of
  translation invariance.
  \choice Average pooling is a linear operation, whereas max pooling is not.
\end{choices}

\begin{solution}
  Answer: C.

  Strides help with dimensionality reduction, and that is all they do. Increasing the stride simply
  samples the convolution at fewer positions --- it discards information (the more so the larger the
  stride) and gives no invariance whatsoever, because shifting the input by one pixel changes which
  positions get sampled. Max pooling is different precisely because of the maximum: the largest
  activation inside the window can move around inside that window and the output does not change.
  That is a real, if partial and purely local, invariance.
\end{solution}

% =============================================================== Question 8 ===
\question
Three standard measures are used to judge how well a statistical shape model has embedded the
training data into a lower-dimensional space. Which of the following statements is
\textbf{incorrect}?

\begin{choices}
  \choice Compactness measures how much of the total variance is explained by the first $t$ modes,
  $C(t) = \sum_{i \leq t} \lambda_i / \lambda_{total}$; a compact model needs few parameters to
  generate new shape instances.
  \choice Generalisation measures the ability of the model to represent instances that were not in
  the training set, and is estimated with leave-one-out reconstruction.
  \choice Specificity measures how accurately the model reconstructs the training shapes
  themselves; a perfectly specific model reproduces every training instance exactly.
  \choice Principal component analysis assumes that the data follow a Gaussian distribution, and
  produces an orthogonal lower-dimensional coordinate system that maximises the variance along each
  successive axis.
\end{choices}

\begin{solution}
  Answer: C.

  Specificity asks the opposite question: does the model generate \emph{only valid} shapes? It is
  measured by drawing a large number of random instances from the model and, for each, computing
  the distance to the closest shape in the training set. A model that can generate anatomically
  impossible shapes scores badly, however well it reconstructs the training data.

  Note the tension this creates with (B): generalisation improves as you keep more modes,
  specificity degrades. The number of modes retained is a trade-off between the two.
\end{solution}

% =============================================================== Question 9 ===
\question
Which of the following statements about positional encoding and the structure of a transformer
layer is \textbf{incorrect}?

\begin{choices}
  \choice Without positional encoding a transformer is equivariant to permutations of its input
  tokens, so ``I bought an apple watch'' and ``watch an apple I bought'' cannot be distinguished.
  \choice An ideal positional encoding should give a unique representation for each position, be
  bounded, generalise to sequences longer than those seen in training, and encode the relative
  positions of tokens.
  \choice In the sinusoidal scheme the positional vector is concatenated to the token embedding,
  which doubles the dimension of the representation entering the first transformer layer.
  \choice A transformer layer consists of multi-head self-attention followed by a per-token MLP,
  each of the two wrapped in a residual connection and followed by layer normalisation.
\end{choices}

\begin{solution}
  Answer: C.

  The positional vector is \textbf{added}, not concatenated: $\tilde{x}_i = x_i + r_i$ with
  $r_i, x_i \in \mathbb{R}^{D \times 1}$. The dimension of the representation is therefore
  unchanged. This is why the encoding has to be bounded --- an unbounded positional vector added to
  the embedding would swamp the token's own content.
\end{solution}

% ============================================================== Question 10 ===
\question
The magnitude image of an MR acquisition is formed from the real and imaginary channels as
$J(x) = \sqrt{(I_r(x) + \epsilon)^2 + (I_i(x) + \eta)^2}$, where $\epsilon, \eta \sim
\mathcal{N}(0, \sigma)$ are the noise in the two channels. What kind of noise does the resulting
magnitude image have, and why does it matter?

\begin{choices}
  \choice Additive Gaussian noise, since $\epsilon$ and $\eta$ are both Gaussian and the operations
  applied to them are continuous.
  \choice Multiplicative speckle noise, of the same type as in ultrasound and optical coherence
  tomography.
  \choice Rician noise --- taking the magnitude makes the noise non-Gaussian and signal-dependent,
  so denoising methods derived from an additive Gaussian model are mis-specified, particularly in
  low-signal regions.
  \choice Poisson noise, of the same type as the quantum noise in low-dose CT.
  \choice No noise at all, because the two independent noise terms cancel in the magnitude
  operation.
\end{choices}

\begin{solution}
  Answer: C.

  The noise is Gaussian in the real and imaginary channels, but the magnitude operation is
  non-linear, and a non-linear function of Gaussian variables is not Gaussian. The result is the
  \textbf{Rician} distribution. The practical consequence is that the noise is signal-dependent:
  where the true signal is strong the Rician distribution is close to Gaussian and an additive
  model is a fair approximation, but in dark, low-signal regions it is strongly skewed and even
  biases the mean intensity upwards. This is why specialised Rician denoisers exist for MRI.
\end{solution}

% ============================================================== Question 11 ===
\question
Consider a classification network that takes grey-scale input images of size $60 \times 60$ with
one input channel. All convolutions are \emph{valid} (no padding, $P = 0$) with stride
$S = 1$, and each is followed by max pooling with $K = 2 \times 2$ and $S = 2$:

\begin{itemize}
  \item \textbf{Conv1}: 12 kernels with $K = 5 \times 5$ \quad\textbf{Pool1}: $2\times2$, $S = 2$
  \item \textbf{Conv2}: 24 kernels with $K = 5 \times 5$ \quad\textbf{Pool2}: $2\times2$, $S = 2$
  \item \textbf{Conv3}: 48 kernels with $K = 3 \times 3$ \quad\textbf{Pool3}: $2\times2$, $S = 2$
\end{itemize}

\noindent Recall that a valid convolution gives $N_l = N_{l-1} - k_1 + 1$. Circle True or False:

\begin{parts}
  \part \TFT\ \ The output of Pool3 has a spatial size of $5 \times 5$.

  \part \TFF\ \ Each neuron in Pool3 has a global receptive field covering the entire
  $60 \times 60$ input image.

  \part \TFT\ \ The total number of convolutional kernel weights (excluding the biases) is
  $5 \cdot 5 \cdot 1 \cdot 12 + 5 \cdot 5 \cdot 12 \cdot 24 + 3 \cdot 3 \cdot 24 \cdot 48
  = 17\,868$.

  \part \TFF\ \ Replacing the three max-pooling layers by $2 \times 2$ convolutions with stride 2
  would leave the total number of learnable parameters unchanged, because pooling and strided
  convolution both simply downsample.
\end{parts}

\begin{solution}
  (a) TRUE --- tracking the sizes: $60 \xrightarrow{\text{Conv1}} 56 \xrightarrow{\text{Pool1}} 28
  \xrightarrow{\text{Conv2}} 24 \xrightarrow{\text{Pool2}} 12 \xrightarrow{\text{Conv3}} 10
  \xrightarrow{\text{Pool3}} 5$.

  (b) FALSE --- it is $28 \times 28$, less than a quarter of the image area. Track the receptive
  field $R$ and jump $j$ from $R = 1$, $j = 1$ with $R \leftarrow R + (K-1)\cdot j$ and
  $j \leftarrow j \cdot S$: Conv1 $\rightarrow R = 5,\ j = 1$; Pool1 $\rightarrow R = 6,\ j = 2$;
  Conv2 $\rightarrow R = 14,\ j = 2$; Pool2 $\rightarrow R = 16,\ j = 4$;
  Conv3 $\rightarrow R = 24,\ j = 4$; Pool3 $\rightarrow R = 28,\ j = 8$. Output size and receptive
  field are different things --- a $5 \times 5$ output does not mean each of those 25 neurons sees
  everything. (The global stride, incidentally, is $8 \times 8$.)

  (c) TRUE --- $300 + 7\,200 + 10\,368 = 17\,868$. Pooling layers contribute nothing.

  (d) FALSE --- max pooling has \emph{no} learnable parameters, whereas convolutions do. The
  replacement would add
  $2 \cdot 2 \cdot 12 \cdot 12 + 2 \cdot 2 \cdot 24 \cdot 24 + 2 \cdot 2 \cdot 48 \cdot 48
  = 576 + 2\,304 + 9\,216 = 12\,096$ new weights, an increase of about 68\%. The two operations
  also differ in behaviour, not just in cost: pooling buys partial local translation invariance,
  strided convolution does not.
\end{solution}

% ============================================================== Question 12 ===
\question
Which of the following statements on domain shift, adaptation and uncertainty are correct (multiple
answers are possible)?

\begin{choices}
  \choice Acquisition shift is a change in $P_D(X \mid Z)$ --- scanner, resolution, contrast,
  protocol --- and is the kind of dataset shift encountered most often in medical imaging.
  \choice Prevalence shift is a change in $P_D(Y \mid X)$, caused for instance by differences in
  annotation policy or in the experience of the annotator.
  \choice The Expected Calibration Error measures the average gap between the confidence a model
  assigns to its predictions and the accuracy it actually achieves at that confidence.
  \choice Adapting only the batch-normalisation parameters of a network is not possible, because
  batch normalisation has no learnable parameters.
  \choice A deep ensemble approximates $p(\theta \mid D, M)$ by training the same architecture
  several times from different random initialisations, on the assumption that the different runs
  land in different modes of the parameter posterior.
\end{choices}

\begin{solution}
  Answer: (A), (C) and (E) are correct.

  (B) confuses two rows of the taxonomy. \emph{Prevalence} shift is a change in $P_D(Y)$ --- the
  case--control balance, or how the target population was selected. A change in $P_D(Y \mid X)$
  driven by annotation policy or annotator experience is \emph{annotation} shift.

  (D) is wrong: batch normalisation computes
  $\mathrm{BN}(a_l) = \gamma\,(a_l - \mu_l)/\sqrt{\sigma_l^2 + \epsilon} + \beta$, and the scale
  $\gamma$ and shift $\beta$ are learnable. Adapting only those is a standard lightweight strategy,
  attractive because it touches very few parameters and leaves the expensively trained weights
  alone --- though it is not obvious that they are the right parameters for correcting a change in
  intensity characteristics.
\end{solution}

\end{questions}

\end{document}
"
%% Compile the blank version (no answers):  pdflatex Sample_Exam_4.tex
%%
\ifx\withanswers\undefined
  \documentclass[11pt]{exam}
\else
  \documentclass[11pt,answers]{exam}
\fi

\usepackage[margin=2.4cm]{geometry}
\usepackage{amsmath,amssymb}
\usepackage{array}
\usepackage[T1]{fontenc}

% ---------------------------------------------------------------- page style
\pagestyle{foot}
\firstpageheader{}{}{}
\firstpagefooter{}{}{}
\runningheader{}{}{}
\runningfooter{}{}{Page \thepage}

% -------------------------------------------------------------- multiplechoice
% Reproduce the "(A) (B) (C) ..." labelling of the original exam.
\renewcommand{\choicelabel}{$\bigcirc$~(\Alph{choice})}

% ---------------------------------------------------------- true/false macros
% \TFT -> the statement is TRUE, \TFF -> the statement is FALSE.
% In the answer version the correct word is set in capitals and bold.
\newcommand{\TFT}{\ifprintanswers\textbf{TRUE}\hspace{1.1em}False\else True\hspace{1.1em}False\fi}
\newcommand{\TFF}{True\hspace{1.1em}\ifprintanswers\textbf{FALSE}\else False\fi}

% ------------------------------------------------------------- writing space
% Reserves room to write in the blank version; collapses in the answer version.
\newcommand{\answerspace}[1]{\ifprintanswers\else\vspace{#1}\fi}

\begin{document}

% ------------------------------------------------------------------ title page
\begin{center}
  {\LARGE\bfseries Sample Exam}\\[2mm]
  {\Large Medical Image Analysis}\\[6mm]
\end{center}

\vspace{4mm}
\noindent Student name:\hrulefill

\vspace{6mm}
\noindent Immatriculation number:\hrulefill

\vspace{10mm}

\begin{center}
  {\large\bfseries Exam Questions}
\end{center}

\noindent\textbf{Note:} All the questions have equal weight.

\vspace{4mm}

\begin{questions}

% =============================================================== Question 1 ===
\question
In 2D acquisitions the slices are acquired one after the other in a pre-defined direction and
stacked into a 3D volume; in axial acquisitions the slices lie in the transverse plane. Voxel
sizes are given as in-plane (two values) and through-plane (one value, the slice thickness), and
the field-of-view is image size $\times$ voxel size. A \emph{dynamic} (4D) acquisition repeats the
whole volume a number of times, one volume per time point.

\vspace{1mm}
\noindent An axial dynamic acquisition is performed with in-plane voxel size
$1.5 \times 1.5$~mm$^2$ and 128 pixels in both directions, through-plane voxel size $5.0$~mm with
30 slices, and 200 time points acquired at a temporal resolution of $2.5$~s per volume.

\begin{parts}

  \part What is the field-of-view in the head-to-toe, left-to-right and anterior-to-posterior
  directions, respectively?

  \begin{solution}
    Answer: $150 \times 192 \times 192$~mm$^3$.

    The through-plane direction of an axial acquisition is head-to-toe: $30 \times 5.0 = 150$~mm.
    Both in-plane directions give $128 \times 1.5 = 192$~mm.
  \end{solution}
\answerspace{2.2cm}

  \part How many voxels does the complete 4D data set contain, and how long does the acquisition
  take?

  \begin{solution}
    Answer: $98\,304\,000$ voxels, and $500$~s $= 8$~min $20$~s.

    Per volume: $128 \times 128 \times 30 = 491\,520$ voxels. Over 200 time points:
    $491\,520 \times 200 = 98\,304\,000$. Duration: $200 \times 2.5 = 500$~s.
  \end{solution}
\answerspace{2.8cm}

  \part A second scan of the same patient is acquired at $1.0 \times 1.0 \times 1.0$~mm$^3$
  isotropic. To compare measurements between the two scans, which property of the images has to
  be matched? Give the resize factors needed to bring one volume of the dynamic acquisition onto
  that grid, and the resulting image size.

  \begin{solution}
    Answer: the \emph{voxel size} has to be matched (never the image size). Resize factors
    $5.0 \times 1.5 \times 1.5$ along head-to-toe, left-to-right and anterior-to-posterior,
    giving $150 \times 192 \times 192$~voxels.

    The resize factor is the ratio of the old voxel size to the new one: $5.0/1.0 = 5.0$
    head-to-toe and $1.5/1.0 = 1.5$ in the two in-plane directions. Applying these:
    $30 \cdot 5.0 = 150$ and $128 \cdot 1.5 = 192$ twice. The field-of-view is unchanged, as it
    must be --- only the sampling of it has changed.
  \end{solution}
\answerspace{4.2cm}

\end{parts}

% =============================================================== Question 2 ===
\question
Otsu's method selects a threshold for foreground--background separation automatically from the
image histogram. Let $p_1, p_2$ be the numbers of pixels assigned to the two groups by a candidate
threshold and $\sigma_1^2, \sigma_2^2$ the intensity variances within those groups. Which of the
two criteria below is the one Otsu's method uses? Please circle the correct choice.

\vspace{2mm}
\begin{center}
  $\displaystyle \min_\tau\ \ p_1 \sigma_1^2 + p_2 \sigma_2^2$
  \hspace{3.2cm}
  $\displaystyle \max_\tau\ \ p_1 \sigma_1^2 + p_2 \sigma_2^2$
\end{center}
\vspace{2mm}

\begin{solution}
  Answer: Left one --- Otsu's method \emph{minimises} the intra-class variance.

  A good threshold is one for which the two resulting groups are each as homogeneous as possible,
  i.e.\ the spread of intensities \emph{within} a group is small. Maximising the same quantity
  would deliberately choose the threshold that makes the two groups as internally inconsistent as
  possible, which is the worst threshold available.
\end{solution}
\answerspace{1.5cm}

% =============================================================== Question 3 ===
\question
Denoising can be written either as an energy minimisation or as a maximum-a-posteriori estimate.
$J$ is the observed noisy image and $I$ the image we want to recover.

\vspace{2mm}
\begin{center}
  $\displaystyle \max_I\ \log P(J \mid I) + \log P(I)$
  \hspace{2.0cm}
  $\displaystyle \min_I\ \lambda\, D(I,J) + R(I)$
\end{center}
\vspace{2mm}

\begin{parts}
  \part Please indicate which term corresponds to which term in the two formulations.

  \begin{solution}
    Answer:
    \begin{align*}
      \log P(J \mid I) &\ \Longleftrightarrow\ \lambda\, D(I,J) &&\text{(data consistency)}\\[1mm]
      \log P(I)        &\ \Longleftrightarrow\ R(I)             &&\text{(regularisation)}
    \end{align*}

    The likelihood encodes how the noise process works; the prior encodes which images are
    considered plausible before any data is seen.
  \end{solution}
\answerspace{2.4cm}

  \part For total variation denoising the data term is $D(I,J) = \|I - J\|_2^2$ and the
  corresponding likelihood is Gaussian, $P(J \mid I) = \mathcal{N}(J;\, I,\, \sigma^2)$. What is
  $\lambda$ in terms of $\sigma$, and what does that mean in words?

  \begin{solution}
    Answer: $\lambda = 1/(2\sigma^2)$.

    Taking $-\log$ of a Gaussian likelihood gives $\|I-J\|_2^2 / (2\sigma^2)$ up to a constant, so
    the weight in front of the data term \emph{is} the inverse noise variance. In words: the
    weighting between data term and regulariser is not an arbitrary tuning knob --- it states how
    much noise you believe is in the observation. A noisier image (large $\sigma$) means a small
    $\lambda$, which trusts the data less and lets the prior smooth more.
  \end{solution}
\answerspace{3.0cm}
\end{parts}

% =============================================================== Question 4 ===
\question
Please indicate whether the statement is true or false by circling the correct choice in each of
the following. $A$ denotes a binary image viewed as a set of pixels, $B$ a structural element, and
$B_x$ the structural element translated to $x$.

\begin{parts}
  \part \TFT\ \ The dilation of $A$ by $B$ is the set of all $x$ for which $B_x$ has a non-empty
  intersection with $A$.

  \part \TFF\ \ The opening of $A$ by $B$ is a dilation followed by an erosion, and it fills small
  gaps and holes.

  \part \TFF\ \ Morphological operations are linear and shift-invariant.

  \part \TFT\ \ Applying a median filter to a label map cannot introduce labels that were not
  already present in it.
\end{parts}

\begin{solution}
  (a) TRUE --- $A \oplus B \triangleq \{x \mid B_x \cap A \neq \emptyset\}$. Erosion is the
  stricter condition $A \ominus B \triangleq \{x \mid B_x \subseteq A\}$.

  (b) FALSE --- both halves are wrong, and they are wrong together. \emph{Opening} is an erosion
  followed by a dilation, $(A \ominus B) \oplus B$, and it removes small islands and protrusions.
  \emph{Closing} is the other order, $(A \oplus B) \ominus B$, and that is the one that fills small
  gaps and holes.

  (c) FALSE --- they are shift-invariant, but they are emphatically \emph{non-linear}. That is
  precisely why they can remove an isolated island without blurring the boundary of the structure
  next to it, which no linear filter can do.

  (d) TRUE --- the median of a set of values is always one of those values, so no new value can
  appear. This is what makes median filtering a safe post-processing step on a segmentation: it
  removes isolated mislabelled voxels without inventing a label that does not exist. Averaging
  filters do not have this property.
\end{solution}

% =============================================================== Question 5 ===
\question
In landmark-based registration a set of corresponding landmark pairs is used to determine the
transformation parameters. How many corresponding landmark pairs are needed for the system to be
exactly determined when the transformation model is (i) a 2D affine transformation, and (ii) a 3D
affine transformation? Explain the counting.

\begin{solution}
  Answer: (i) \textbf{3} landmark pairs in 2D, (ii) \textbf{4} landmark pairs in 3D.

  A 2D affine transformation has 6 free parameters ($A \in \mathbb{R}^{2\times2}$ plus
  $t \in \mathbb{R}^2$). Each landmark correspondence must match both coordinates of the point and
  therefore contributes 2 scalar equations. Setting $2N = 6$ gives $N = 3$.

  A 3D affine transformation has 12 free parameters ($A \in \mathbb{R}^{3\times3}$ plus
  $t \in \mathbb{R}^3$), and each correspondence contributes 3 equations, so $3N = 12$ gives
  $N = 4$.

  With more landmarks than this the system becomes over-determined and is solved in the
  least-squares sense: the landmarks themselves are then no longer aligned exactly, but the rest of
  the image is aligned better --- which is usually what you actually want.
\end{solution}
\answerspace{4.0cm}

% =============================================================== Question 6 ===
\question
Optimisation in image registration is difficult, and the objective function generally has local
minima that a gradient-based optimiser can get stuck in. Please describe the two practical
strategies presented in the lecture for dealing with this, and briefly explain why each of them
helps. You do not need to write down the equations.

\begin{solution}
  Answer: \emph{Sequential optimisation.} First align the images with a \textbf{linear}
  transformation model (rigid or affine), and only then run the deformable, non-linear
  registration. The linear model has very few parameters and removes the large global
  misalignment, so the non-linear stage starts close to the correct answer and only has to
  explain the local differences that remain. Running the deformable model straight away would ask
  it to absorb a global rotation or translation through local deformation, which is both harder to
  optimise and physically wrong.

  \emph{Multi-scale optimisation.} Build an image pyramid and start the registration at the
  \textbf{coarsest} scale, using the result of each scale to initialise the next finer one,
  $\theta^{(k-1)}_0 = \theta^{(k)}_*$. A coarse image has fewer details, so its objective function
  is smoother and has fewer local minima; the coarse solution captures the large-scale alignment
  and places the finer optimisation inside the basin of the correct minimum.

  The two are routinely combined, and both are really the same idea: solve an easier version of the
  problem first and use its answer as the starting point for the harder one.
\end{solution}
\answerspace{7.5cm}

% =============================================================== Question 7 ===
\question
Convolutional networks reduce the spatial dimensions of their feature maps either with pooling or
with strided convolutions. Which of the following statements is \textbf{inaccurate}?

\begin{choices}
  \choice Max pooling provides partial local translation invariance: the position of the maximum
  can move within the pooling window without changing the pooled value.
  \choice Pooling is applied to each channel separately and has no learnable parameters at all.
  \choice A strided convolution reduces the spatial dimensions and, like pooling, buys a degree of
  translation invariance.
  \choice Average pooling is a linear operation, whereas max pooling is not.
\end{choices}

\begin{solution}
  Answer: C.

  Strides help with dimensionality reduction, and that is all they do. Increasing the stride simply
  samples the convolution at fewer positions --- it discards information (the more so the larger the
  stride) and gives no invariance whatsoever, because shifting the input by one pixel changes which
  positions get sampled. Max pooling is different precisely because of the maximum: the largest
  activation inside the window can move around inside that window and the output does not change.
  That is a real, if partial and purely local, invariance.
\end{solution}

% =============================================================== Question 8 ===
\question
Three standard measures are used to judge how well a statistical shape model has embedded the
training data into a lower-dimensional space. Which of the following statements is
\textbf{incorrect}?

\begin{choices}
  \choice Compactness measures how much of the total variance is explained by the first $t$ modes,
  $C(t) = \sum_{i \leq t} \lambda_i / \lambda_{total}$; a compact model needs few parameters to
  generate new shape instances.
  \choice Generalisation measures the ability of the model to represent instances that were not in
  the training set, and is estimated with leave-one-out reconstruction.
  \choice Specificity measures how accurately the model reconstructs the training shapes
  themselves; a perfectly specific model reproduces every training instance exactly.
  \choice Principal component analysis assumes that the data follow a Gaussian distribution, and
  produces an orthogonal lower-dimensional coordinate system that maximises the variance along each
  successive axis.
\end{choices}

\begin{solution}
  Answer: C.

  Specificity asks the opposite question: does the model generate \emph{only valid} shapes? It is
  measured by drawing a large number of random instances from the model and, for each, computing
  the distance to the closest shape in the training set. A model that can generate anatomically
  impossible shapes scores badly, however well it reconstructs the training data.

  Note the tension this creates with (B): generalisation improves as you keep more modes,
  specificity degrades. The number of modes retained is a trade-off between the two.
\end{solution}

% =============================================================== Question 9 ===
\question
Which of the following statements about positional encoding and the structure of a transformer
layer is \textbf{incorrect}?

\begin{choices}
  \choice Without positional encoding a transformer is equivariant to permutations of its input
  tokens, so ``I bought an apple watch'' and ``watch an apple I bought'' cannot be distinguished.
  \choice An ideal positional encoding should give a unique representation for each position, be
  bounded, generalise to sequences longer than those seen in training, and encode the relative
  positions of tokens.
  \choice In the sinusoidal scheme the positional vector is concatenated to the token embedding,
  which doubles the dimension of the representation entering the first transformer layer.
  \choice A transformer layer consists of multi-head self-attention followed by a per-token MLP,
  each of the two wrapped in a residual connection and followed by layer normalisation.
\end{choices}

\begin{solution}
  Answer: C.

  The positional vector is \textbf{added}, not concatenated: $\tilde{x}_i = x_i + r_i$ with
  $r_i, x_i \in \mathbb{R}^{D \times 1}$. The dimension of the representation is therefore
  unchanged. This is why the encoding has to be bounded --- an unbounded positional vector added to
  the embedding would swamp the token's own content.
\end{solution}

% ============================================================== Question 10 ===
\question
The magnitude image of an MR acquisition is formed from the real and imaginary channels as
$J(x) = \sqrt{(I_r(x) + \epsilon)^2 + (I_i(x) + \eta)^2}$, where $\epsilon, \eta \sim
\mathcal{N}(0, \sigma)$ are the noise in the two channels. What kind of noise does the resulting
magnitude image have, and why does it matter?

\begin{choices}
  \choice Additive Gaussian noise, since $\epsilon$ and $\eta$ are both Gaussian and the operations
  applied to them are continuous.
  \choice Multiplicative speckle noise, of the same type as in ultrasound and optical coherence
  tomography.
  \choice Rician noise --- taking the magnitude makes the noise non-Gaussian and signal-dependent,
  so denoising methods derived from an additive Gaussian model are mis-specified, particularly in
  low-signal regions.
  \choice Poisson noise, of the same type as the quantum noise in low-dose CT.
  \choice No noise at all, because the two independent noise terms cancel in the magnitude
  operation.
\end{choices}

\begin{solution}
  Answer: C.

  The noise is Gaussian in the real and imaginary channels, but the magnitude operation is
  non-linear, and a non-linear function of Gaussian variables is not Gaussian. The result is the
  \textbf{Rician} distribution. The practical consequence is that the noise is signal-dependent:
  where the true signal is strong the Rician distribution is close to Gaussian and an additive
  model is a fair approximation, but in dark, low-signal regions it is strongly skewed and even
  biases the mean intensity upwards. This is why specialised Rician denoisers exist for MRI.
\end{solution}

% ============================================================== Question 11 ===
\question
Consider a classification network that takes grey-scale input images of size $60 \times 60$ with
one input channel. All convolutions are \emph{valid} (no padding, $P = 0$) with stride
$S = 1$, and each is followed by max pooling with $K = 2 \times 2$ and $S = 2$:

\begin{itemize}
  \item \textbf{Conv1}: 12 kernels with $K = 5 \times 5$ \quad\textbf{Pool1}: $2\times2$, $S = 2$
  \item \textbf{Conv2}: 24 kernels with $K = 5 \times 5$ \quad\textbf{Pool2}: $2\times2$, $S = 2$
  \item \textbf{Conv3}: 48 kernels with $K = 3 \times 3$ \quad\textbf{Pool3}: $2\times2$, $S = 2$
\end{itemize}

\noindent Recall that a valid convolution gives $N_l = N_{l-1} - k_1 + 1$. Circle True or False:

\begin{parts}
  \part \TFT\ \ The output of Pool3 has a spatial size of $5 \times 5$.

  \part \TFF\ \ Each neuron in Pool3 has a global receptive field covering the entire
  $60 \times 60$ input image.

  \part \TFT\ \ The total number of convolutional kernel weights (excluding the biases) is
  $5 \cdot 5 \cdot 1 \cdot 12 + 5 \cdot 5 \cdot 12 \cdot 24 + 3 \cdot 3 \cdot 24 \cdot 48
  = 17\,868$.

  \part \TFF\ \ Replacing the three max-pooling layers by $2 \times 2$ convolutions with stride 2
  would leave the total number of learnable parameters unchanged, because pooling and strided
  convolution both simply downsample.
\end{parts}

\begin{solution}
  (a) TRUE --- tracking the sizes: $60 \xrightarrow{\text{Conv1}} 56 \xrightarrow{\text{Pool1}} 28
  \xrightarrow{\text{Conv2}} 24 \xrightarrow{\text{Pool2}} 12 \xrightarrow{\text{Conv3}} 10
  \xrightarrow{\text{Pool3}} 5$.

  (b) FALSE --- it is $28 \times 28$, less than a quarter of the image area. Track the receptive
  field $R$ and jump $j$ from $R = 1$, $j = 1$ with $R \leftarrow R + (K-1)\cdot j$ and
  $j \leftarrow j \cdot S$: Conv1 $\rightarrow R = 5,\ j = 1$; Pool1 $\rightarrow R = 6,\ j = 2$;
  Conv2 $\rightarrow R = 14,\ j = 2$; Pool2 $\rightarrow R = 16,\ j = 4$;
  Conv3 $\rightarrow R = 24,\ j = 4$; Pool3 $\rightarrow R = 28,\ j = 8$. Output size and receptive
  field are different things --- a $5 \times 5$ output does not mean each of those 25 neurons sees
  everything. (The global stride, incidentally, is $8 \times 8$.)

  (c) TRUE --- $300 + 7\,200 + 10\,368 = 17\,868$. Pooling layers contribute nothing.

  (d) FALSE --- max pooling has \emph{no} learnable parameters, whereas convolutions do. The
  replacement would add
  $2 \cdot 2 \cdot 12 \cdot 12 + 2 \cdot 2 \cdot 24 \cdot 24 + 2 \cdot 2 \cdot 48 \cdot 48
  = 576 + 2\,304 + 9\,216 = 12\,096$ new weights, an increase of about 68\%. The two operations
  also differ in behaviour, not just in cost: pooling buys partial local translation invariance,
  strided convolution does not.
\end{solution}

% ============================================================== Question 12 ===
\question
Which of the following statements on domain shift, adaptation and uncertainty are correct (multiple
answers are possible)?

\begin{choices}
  \choice Acquisition shift is a change in $P_D(X \mid Z)$ --- scanner, resolution, contrast,
  protocol --- and is the kind of dataset shift encountered most often in medical imaging.
  \choice Prevalence shift is a change in $P_D(Y \mid X)$, caused for instance by differences in
  annotation policy or in the experience of the annotator.
  \choice The Expected Calibration Error measures the average gap between the confidence a model
  assigns to its predictions and the accuracy it actually achieves at that confidence.
  \choice Adapting only the batch-normalisation parameters of a network is not possible, because
  batch normalisation has no learnable parameters.
  \choice A deep ensemble approximates $p(\theta \mid D, M)$ by training the same architecture
  several times from different random initialisations, on the assumption that the different runs
  land in different modes of the parameter posterior.
\end{choices}

\begin{solution}
  Answer: (A), (C) and (E) are correct.

  (B) confuses two rows of the taxonomy. \emph{Prevalence} shift is a change in $P_D(Y)$ --- the
  case--control balance, or how the target population was selected. A change in $P_D(Y \mid X)$
  driven by annotation policy or annotator experience is \emph{annotation} shift.

  (D) is wrong: batch normalisation computes
  $\mathrm{BN}(a_l) = \gamma\,(a_l - \mu_l)/\sqrt{\sigma_l^2 + \epsilon} + \beta$, and the scale
  $\gamma$ and shift $\beta$ are learnable. Adapting only those is a standard lightweight strategy,
  attractive because it touches very few parameters and leaves the expensively trained weights
  alone --- though it is not obvious that they are the right parameters for correcting a change in
  intensity characteristics.
\end{solution}

\end{questions}

\end{document}
